% TOP_PROBLEM: {"id": 3088, "title": "Partitioning the Projective Plane", "category_id": 6, "category": {"id": 6, "name": "geometry", "display_name": "Geometry", "description": "Euclidean and non-Euclidean geometry, geometric structures.", "slug": "geometry", "order_index": 6, "created_at": "2026-07-31T15:26:25.671Z"}, "status": "open", "problem_number": "OPG-56328", "set_id": 12}
% FIRST_POSTED: 2026-10-01
% ATTEMPT_STATUS: unresolved
% COMPLETION: 5
% MODEL: GPT 6 Astra Ultra
% UnsolvedMath ID 3088; source catalogue code OPG-56328
% !TeX program = lualatex
% Research attempt, 2026-10-01. Canonical statement preserved verbatim.
\documentclass[11pt,a4paper]{article}
\usepackage[margin=25mm]{geometry}
\usepackage{fontspec}
\setmainfont{lmroman10-regular.otf}[BoldFont=lmroman10-bold.otf,
 ItalicFont=lmroman10-italic.otf,BoldItalicFont=lmroman10-bolditalic.otf]
\usepackage{amsmath,amssymb,mathtools}
\usepackage{unicode-math}
\setmathfont{latinmodern-math.otf}
\AtBeginDocument{\let\setminus\smallsetminus}
\usepackage{xurl,hyperref}
\hypersetup{colorlinks=true,urlcolor=blue}
\setcounter{secnumdepth}{0}
\setlength{\parindent}{0pt}
\setlength{\parskip}{5pt}
\setlength{\emergencystretch}{3em}
\begin{document}
\section*{Partitioning the Projective Plane}\label{3088}
{\small UnsolvedMath ID 3088 · OPG-56328 · Geometry · OpenGarden}
\subsection{Definitions and mathematical statement}
Write $\mathbb{RP}^2$ for the set of one-dimensional linear subspaces of
$\mathbb R^3$. Its elements are unoriented lines through the origin.
For an orthonormal basis $e=(e_1,e_2,e_3)$ define
\[
 O_e=\{\mathbb R x:x\ne0,\quad
                x=t_1e_1+t_2e_2+t_3e_3,\quad t_1,t_2,t_3\ge0\}.
\]
A set of lines $S$ is octahedral when $S\subseteq O_e$ for some basis $e$.
This agrees with the source's geometric definition: every line meets
the closed face $\operatorname{conv}\{e_1,e_2,e_3\}$ and its opposite in
the regular octahedron with vertices $\pm e_i$. Zero coordinates are allowed,
so face edges and vertices are included.

Call $S$ weakly octahedral if each three-element subset of $S$ is
contained in some $O_e$, where the basis may depend on the triple.
The conjecture states that if
\[
 \mathbb{RP}^2=S_1\mathbin{\dot\cup}S_2\mathbin{\dot\cup}S_3
                         \mathbin{\dot\cup}S_4
\]
and every $S_i$ is weakly octahedral, then every $S_i$ is octahedral.
Thus for each part there should be one orthonormal basis working for
all its lines simultaneously. Different parts may use different bases.

The partition must cover every projective line, with disjoint parts;
no measurability or closedness assumption is imposed. The local triple
condition alone, for an arbitrary subset not occurring as a part of such
a partition, is not the assertion. Directions are identified up to all
nonzero scalar multiples, so the choice of orientation of representatives
cannot affect either condition.
\subsection{Short English statement}
If the real projective plane is partitioned into four sets whose every triple fits in an octahedral face pair, must each entire set fit in one?
\subsection{Research attempt}
\paragraph{Scope and source check.}
The original posting by Jonathan Noel [S2], checked 1 October 2026, explicitly
warns that the four-part partition hypothesis matters. We investigate the
obstruction by proving the acute-cap example and a compactness reduction.
Neither result assumes that the parts are measurable, and neither purports
to construct a counterexample partition.

\paragraph{Three acute vectors fit in an orthant.}
Let $x,y,z$ be distinct unit representatives with pairwise nonnegative inner
products. Relabel them so that $c=\langle y,z\rangle$ is the largest of the
three products, and put $a=\langle x,y\rangle$, $b=\langle x,z\rangle$.
Since $0\le a,b,c\le1$, we have $c\ge ab$. For distinct unoriented lines,
$a<1$. Choose $e_1=x$, $e_2=(y-ax)/\sqrt{1-a^2}$, and choose the sign of the
remaining perpendicular unit vector $e_3$ so that $\langle z,e_3\rangle\ge0$.
In this orthonormal frame the three coordinate rows are
\[
 (1,0,0),\quad(a,\sqrt{1-a^2},0),\quad
 \left(b,\frac{c-ab}{\sqrt{1-a^2}},\langle z,e_3\rangle\right).
\]
Every entry is nonnegative. Degenerate linear dependence simply makes the
last entry zero. Hence these three lines are octahedral.

\paragraph{An infinite weakly octahedral set with no global frame.}
Fix a unit vector $u$ and let $C$ be the open spherical cap of vectors
$x\in S^2$ with $\langle x,u\rangle>1/\sqrt2$. Let $S$ be their unoriented
lines. Any two vectors in $C$ have angle strictly below $\pi/2$, by adding
their angles from $u$. Thus every triple fits in an orthant by the preceding
calculation, and $S$ is weakly octahedral.
Suppose all its lines belonged to a single $O_e$. The connected cap $C$
would lie in the union of the positive and negative closed orthants for
that frame. Those orthants are disjoint on the unit sphere, so connectedness
puts the whole cap in one of them; reversing all three axes makes it positive.
Each coordinate functional must therefore be nonnegative throughout $C$.
This forces $\langle e_i,u\rangle\ge1/\sqrt2$ for every $i$: otherwise a vector
at angle less than $\pi/4$ from $u$, tilted away from $e_i$, has negative
$i$th coordinate. But orthonormality gives
\[
 1=\|u\|^2=\sum_{i=1}^3\langle e_i,u\rangle^2\ge3/2,
\]
a contradiction. This reproduces, with explicit verification, the obstruction
reported by [S2]; it is not asserted to be a new counterexample.

\paragraph{A finite-witness reduction.}
For each projective line $\ell$ let $F_\ell\subset O(3)$ consist of frames
whose positive or negative orthant contains a unit representative of $\ell$.
It is closed: for a fixed representative it is a union of two sets defined
by three weak continuous inequalities. The space $O(3)$ is compact, since
it is closed and bounded in the space of real $3$ by $3$ matrices.
If every finite subset of $S$ has a common frame, the family
$\{F_\ell:\ell\in S\}$ has the finite intersection property. Compactness
gives a frame for all of $S$. Consequently every non-octahedral set has a
\emph{finite} non-octahedral subset. For a weakly octahedral set that witness
must have at least four lines; sets of at most two lines fit by choosing
signs with nonnegative inner product.

\paragraph{Verification and remaining gap.}
The cap calculation checks both orientation invariance and zero-coordinate
boundary conventions. The compactness argument does not upgrade triplewise
intersection to finite intersection: the cap proves that such an upgrade is
false. A proof of the catalogue conjecture must use the other three parts to
exclude every finite obstruction inside a part. No extension of the cap to
a four-part weakly octahedral partition has been found here.

\paragraph{Final status.}
\textbf{UNRESOLVED.} The output is a proved local obstruction and a finite-witness
reduction, not a proof or disproof of the partition assertion. Completion
estimate: 5\%.

\subsection{Sources}
\begin{itemize}
\item[S1] ULAM.AI and the UnsolvedMath Contributors, supplied problems.json, record 3088 (OPG-56328), OpenGarden collection. Catalogue identity and source transcription; CC BY 4.0.
\url{https://huggingface.co/datasets/ulamai/UnsolvedMath}.
\item[S2] Open Problem Garden, Partitioning the Projective Plane. Original problem and discussion; the mathematical formulation below is expanded from the supplied source record.
\url{http://www.openproblemgarden.org/op/partitioning_the_projective_plane}.

\end{itemize}
\end{document}
